\documentclass[10pt,a4paper]{amsart}
\usepackage[margin=19mm]{geometry}
\usepackage{amsmath,amssymb,mathtools}
\usepackage[scr=rsfs,cal=euler]{mathalfa}
\usepackage{xfrac}
\usepackage[unicode,hidelinks,bookmarks=false]{hyperref}
\newcommand{\ie}{i.e.\ }
\newcommand{\cf}{cf.\ }
\newcommand{\resp}{respectively\ }
\newcommand{\Subsection}{\subsection}
\providecommand{\nobreakdash}{\nobreak\hbox{-}}
\setlength{\emergencystretch}{2em}
\multlinegap0pt
\usepackage{bm}
\usepackage{bbm}
\usepackage{dsfont}
\usepackage{xspace}
\usepackage{cancel}
\usepackage{mathtools}
\usepackage{amsthm}
\makeatletter
\@ifundefined{theorem}{\newtheorem{theorem}{Theorem}}{}
\@ifundefined{lemma}{\newtheorem{lemma}[theorem]{Lemma}}{}
\@ifundefined{remark}{\newtheorem{remark}[theorem]{Remark}}{}
\makeatother
\newcommand{\M}{\mathcal{M}}
\newcommand{\Rn}{\mathbb{R}^{n+1}}
\newcommand{\Sn}{\mathbb{S}^n}
\title[Uniqueness of solutions to a class of isotropic curvature pr]{Uniqueness of solutions to a class of isotropic curvature problems}
\author{Mohammad N. Ivaki, Emanuel Milman}
\date{Source: arXiv:2304.12839v1 (CC BY 4.0)}
\begin{document}
\maketitle

\noindent\emph{Source and licence.} Uniqueness of solutions to a class of isotropic curvature problems. Version arXiv:2304.12839v1 (CC BY 4.0), distributed under the Creative Commons Attribution 4.0 International licence, as recorded in the arXiv licence metadata for this exact version and shown on the abstract page https://arxiv.org/abs/2304.12839v1. Full rights notice: licenses/arxiv-2304.12839-CC-BY-4.0.txt.\par\smallskip

\noindent\textit{Editorial scope.} Counted unit: the Theorem whose source label is \texttt{thm:BCD} in the article (source0.tex lines 290--292, source label \texttt{thm:BCD}), delivered together with the article's own complete original proof of it (source0.tex lines 293--307, one \texttt{proof} environment of 15 lines). No mathematics is written, repaired or re-derived: statement and proof are the article's own text line for line. Required context retained uncounted: magic inequality (lines 204--217); magic ineq statement (lines 210--212); affine ver (lines 214--216); centroid (lines 242--249). Every definition, display and label of those lines is retained unchanged. Omitted from this package and not redistributed: 1--203, 213--213, 217--241, 250--289, 308--458. Nothing inside a retained excerpt is omitted or altered: every retained body is delivered exactly as the article's lines are written, so no omission receipt of any kind accompanies this package. Citations: the retained excerpts carry no citation command at all, so no bibliography entry of the article is transcribed and no reference is redistributed. Reference adaptation: none was needed and none was made; every reference inside the delivered unit resolves inside it, so no unresolved reference or citation is produced. Printed slips kept literal: no printed word, symbol, hypothesis or number of the counted statement or of its proof was altered. Layout and class: the article's own class is \texttt{amsart} with packages including amsrefs, amssymb, amsmath, amsthm, verbatim, xcolor, hyperref, lipsum, cleveref, bookmark, thmtools, mathtools; the wrapper is the lane's pinned \texttt{amsart} editorial wrapper at 10pt on A4, which restates the theorem-like environments the retained text needs and adds the article's own macro definitions that the retained text needs (\texttt{\textbackslash M}, \texttt{\textbackslash Rn}, \texttt{\textbackslash Sn}), with extra wrapper packages (bm, bbm, dsfont, xspace, cancel, mathtools). The delivered numbering is the unit's own. Assets: no retained span contains a figure environment, a graphics-inclusion command, a PDF-page inclusion, a table or a TikZ picture; no image file is copied into this package, the package declares no asset dependency and no image file is redistributed with it. Licence: version arXiv:2304.12839v1 of this article is distributed under the Creative Commons Attribution 4.0 International licence, as recorded in the arXiv licence metadata for this exact version and shown on the abstract page https://arxiv.org/abs/2304.12839v1. A full rights notice is delivered at licenses/arxiv-2304.12839-CC-BY-4.0.txt. Characters: every retained excerpt is pure ASCII, so the delivered engine, XeLaTeX with its default Latin Modern text font, reports no missing character.
\begin{lemma}[Main Lemma]\label{magic inequality}
Let $X=Dh: \Sn\to \partial K$. Then we have
\begin{align*}
k\int |X|^2 dV_k\leq \int h \left ( \sigma_1-(k+1)\frac{\sigma_{k+1}}{\sigma_k} \right )dV_k +k\frac{|\int X dV_k|^2}{\int dV_k}.
\end{align*}
In particular, for $k=n$ we have
\begin{align}\label{magic ineq statement}
n\int|X|^2dV\leq \int h (\bar{\Delta}h+nh)dV+n\frac{\left|\int X dV\right|^2}{\int dV} .
\end{align}
Equivalently, for $k=n$ there holds
\begin{align}\label{affine ver}
\int \langle hX, \bar{\nabla}\log \frac{h^{n+2}}{\mathcal{K}} \rangle dV = \int (n|\bar\nabla h|_{\bar g}^2 - h \bar \Delta h) dV \leq n\frac{|\int XdV|^2}{\int dV}.
\end{align}
\end{lemma}

\begin{remark}\label{centroid}
By the divergence theorem, for any vector $w \in \mathbb{R}^{n+1}$,
\begin{align*}
\int \langle X,w\rangle dV & =\int_{\partial K} \langle p,w\rangle\langle p,\nu(p)\rangle \mathcal{H}^n(dp) \\
& =\int_{K} \operatorname{div}_{\Rn}(\langle x,w\rangle x ) dx =(n+2)\int_K \langle x,w\rangle dx.
\end{align*}
Therefore, the vector $\int XdV$ appearing in \eqref{magic ineq statement} and \eqref{affine ver} is a multiple of the centroid of $K$, and is equal to $0$ whenever $\M^n = \partial K$ is origin-centred. 
\end{remark}

\begin{theorem} \label{thm:BCD}
Let $\M^n$ be a smooth, strictly convex hypersurface. If $\mathcal{K}=h^{1-p}$ with $p > -(n+1)$, and either $\M^n$ is origin-centred or $p \leq -1$, then $\M^n$ is an origin-centred ball. 
\end{theorem}

\begin{proof}
%Next we consider the characterization of solutions to $\mathcal{K}=h^{1-p}$ for $p\in[-(n+1), -1]$. 
When $dV = h^p d\mu$, \autoref{magic inequality} and integration by parts yield
\begin{align} \label{p-inq}
\frac{n+1+p}{n}\int |\bar{\nabla} h|_{\bar{g}}^2dV\leq  \frac{\left|\int XdV\right|^2}{\int dV}.
\end{align}
By \autoref{centroid}, we conclude $h$ is constant when $p > -(n+1)$ and $\M^n$ is origin-centred. In the general case, by the divergence theorem,
\begin{align}\label{integ estim}
\int X dV =\frac{n+1+p}{n}\int h^p\bar{\nabla}h d\mu.
\end{align}
Plugging \eqref{integ estim} into \eqref{p-inq}, we deduce when $p > -(n+1)$:
%For $p>-(n+1)$, due to \eqref{integ estim}, $dV=h^pd\mu$, and \autoref{magic inequality} we obtain
\[\int \left|\bar{\nabla}h-\frac{\int \bar{\nabla}hdV}{\int dV}\right|_{\delta}^2dV\leq \frac{p+1}{n}\frac{|\int \bar{\nabla}hdV|_{\delta}^2}{\int dV}.\]
Consequently, when in addition $p \leq -1$, we conclude that $h$ is constant.
\end{proof}

\end{document}
